\chapterstar{Justificación}

Actualmente Farmacia Selene lleva una gestión tradicional, registrando cada articulo, venta y precio de forma manual, lo cual hace que carezca de organización, perdiendo tiempo al momento de hacer el corte de cada mes. Implementaremos una solución tecnológica que unifique la gestión de inventario y ventas junto con el monitoreo ambiental, aportando beneficios directos a Farmacia Selene, optimiza recursos, reduce tiempos de gestión de inventario, gestión de venta y aumentara la productividad. A diferencia de los sistemas tradicionales, esta propuesta combina disciplinas como la computación, la electrónica, el IoT, las telecomunicaciones y el cómputo en la nube, para ofrecer una plataforma integral, escalable y accesible desde cualquier lugar. La integración de sensores de humedad permitirá la generación de alertas criticas automatizadas en tiempo real, una alteración de temperatura puede comprometer la calidad de medicamentos o alimentos, no solo se registrara, mandara al administrador una notificación instantánea e incluso sugerirá acciones para solucionar dicho problema, reduciendo el tiempo de reacción, asegurando la conservación de los productos y la continuidad de las operaciones. La unión de Python, fastAPI, PostgreSQL, garantiza un desarrollo moderno y alto rendimiento, esto asegura que la plataforma pueda evolucionar con las necesidades futuras del negocio. La arquitectura cloud y el frontend asegura que los administradores puedan acceder con facilidad a la información en tiempo real desde cualquier lugar, mejorando la supervisión y la capacidad de respuesta inmediata. Con el fin de ofrecer a la Farmacia Selene una herramienta confiable, adaptable y orientada a la optimización de sus operaciones.

\thispagestyle{empty}
